%% Question bank chapter file - generated by bank_tools.py build (do not edit while the bank is still being built:
%% it is regenerated from the per-test data in Question Bank\data). Every question is one bankq block:
%%   \begin{bankq}[id=..., source={...}, paper=.., qno=.., type=..., subtopic={...}, answer={...}, answer_src=key|solved,
%%                 working={...}, difficulty=1-10, diff_parts={K S I C F}, diff_note={...},
%%                 group=<passage id or empty>, also_in={duplicate copies in other tests}]
%% difficulty = min(10, K+S+I+C+F-1): K concepts 1-3, S steps 1-4, I insight 0-2, C computation 0-2, F format 0-1
%%   ...question LaTeX, options, figures (\bfig / \bfigi -> img/)...
%%   \end{bankq}
%% Passages shared by several questions are bankpassage blocks; their questions carry group=<passage id>.
\documentclass[10pt]{article}
\usepackage[a4paper,margin=1.8cm]{geometry}
\usepackage[T1]{fontenc}
\usepackage{amsmath,amssymb,mathtools}
\usepackage{graphicx}
\usepackage[export]{adjustbox}
\usepackage{xcolor,needspace,keyval}
\usepackage{enumitem}
\usepackage{booktabs,array,tabularx}
\usepackage[hidelinks]{hyperref}
\setlength{\parindent}{0pt}
\setlength{\parskip}{3pt}
\makeatletter
\define@key{bq}{id}{\def\bq@id{#1}}\define@key{bq}{source}{\def\bq@src{#1}}\define@key{bq}{paper}{\def\bq@paper{#1}}
\define@key{bq}{qno}{\def\bq@qno{#1}}\define@key{bq}{type}{\def\bq@type{#1}}\define@key{bq}{subtopic}{\def\bq@sub{#1}}
\define@key{bq}{answer}{\def\bq@ans{#1}}\define@key{bq}{answer_src}{\def\bq@asrc{#1}}\define@key{bq}{working}{\def\bq@work{#1}}
\define@key{bq}{group}{\def\bq@group{#1}}\define@key{bq}{also_in}{\def\bq@also{#1}}
\define@key{bq}{difficulty}{\def\bq@diff{#1}}\define@key{bq}{diff_parts}{\def\bq@dparts{#1}}\define@key{bq}{diff_note}{\def\bq@dnote{#1}}
\newenvironment{bankq}[1][]{\def\bq@work{}\def\bq@also{}\def\bq@group{}\def\bq@diff{?}\def\bq@dparts{}\def\bq@dnote{}\setkeys{bq}{#1}%
  \par\needspace{4\baselineskip}\medskip\noindent{\small\textbf{\bq@id}\hfill\color{gray}[\bq@src, Paper \bq@paper, Q\bq@qno]
  \bq@type\ $\cdot$ \bq@sub\ $\cdot$ \textbf{D\bq@diff}}\par\nopagebreak\relax}
  {\par{\small\color{teal}Answer: \textbf{\bq@ans}\ifx\bq@work\@empty\else\ (solved: \bq@work)\fi
  \ifx\bq@also\@empty\else\par\color{gray}Also in: \bq@also\fi}\par}
\define@key{bp}{id}{\def\bp@id{#1}}\define@key{bp}{source}{\def\bp@src{#1}}\define@key{bp}{questions}{\def\bp@qs{#1}}
\newenvironment{bankpassage}[1][]{\setkeys{bp}{#1}\par\needspace{4\baselineskip}\bigskip\noindent
  {\small\textbf{Passage \bp@id}\hfill\color{gray}[\bp@src; questions \bp@qs]}\par\nopagebreak\relax\itshape}{\par}
\makeatother
\newcommand{\bfig}[2]{\par\smallskip{\centering\includegraphics[width=#1pt,max width=0.62\linewidth,max height=0.32\textheight]{img/#2}\par}\smallskip}
\newcommand{\bfigi}[2]{\includegraphics[width=#1pt,max width=\linewidth,max height=0.32\textheight,valign=c]{img/#2}}
\newcommand{\bmissing}[1]{\textit{[figure missing in the source conversion: \detokenize{#1}]}}
\newenvironment{qoptions}{\begin{itemize}[leftmargin=1.8em,itemsep=1pt,topsep=2pt,label={}]}{\end{itemize}}
\newcommand{\opt}[2]{\item[(#1)] #2}
\providecommand{\fwstart}{}\providecommand{\fwend}{}\providecommand{\qfullsplit}{}
\newcolumntype{L}[1]{>{\raggedright\arraybackslash}p{#1}}
\newcolumntype{C}[1]{>{\centering\arraybackslash}p{#1}}
\newcolumntype{R}[1]{>{\raggedleft\arraybackslash}p{#1}}
\providecommand{\sectionhdr}[1]{}\providecommand{\subsectionhdr}[1]{}


\begin{document}
{\Large\bfseries Coordination Compounds --- Question Bank}\par\smallskip
{\small Chemistry chapter C31. Questions from all converted test papers, grouped by source series. Answers: printed key where it matches the paper, otherwise solved (answer\_src=solved, with working). Exact repeats across papers appear once, with also\_in listing the other copies.}\par\medskip\hrule

\section*{Score-I papers}
\begin{bankq}[id=SI12-P1-CHEM-S3-Q09, source={Score-I - Test 12 (09 Mar 2022)}, paper=1, qno=9, type=numerical, subtopic={Nomenclature}, answer={2}, answer_src=key, difficulty=5, diff_parts={K2 S3 I0 C0 F1}, diff_note={Identify bonding atoms of ligands to see M-C bonds (olefin pi, CO)}]
Find total number of complex in which metal-carbon linkage(s) is/are present.\\\relax
[Fe(EDTA)]$^-$, K[PtCl$_3$($\pi$-C$_2$H$_4$)], [Cr(acac)$_3$], [Fe(CO)$_5$], [Co(C$_2$H$_5$O)$_6$]$^{3-}$, [Ni(dmg)$_2$]
\end{bankq}
\begin{bankpassage}[id=SI12-P2-CHEM-S1-G3-4, source={Score-I - Test 12 (09 Mar 2022)}, questions={3-4}]
\begin{center}\small
\begin{tabular}{|p{0.8cm}|p{5.0cm}|p{0.8cm}|p{3.8cm}|p{0.8cm}|p{2.8cm}|}
\hline
 & \centering Column-I\newline Compound & & \centering Column-II\newline Hybridisation and magnetic behaviour & & \centering Column-III\newline Isomerism\tabularnewline
\hline
(I) & [Cr(NH$_3$)$_4$Cl$_2$]Cl & (i) & Paramagnetic & (P) & Ionisation\\
(II) & [Ti(H$_2$O)$_5$Cl](NO$_3$)$_2$ & (ii) & Diamagnetic & (Q) & Geometrical\\
(III) & [Pt(en)(NH$_3$)Cl]NO$_3$ & (iii) & d$^2$sp$^3$ & (R) & Linkage\\
(IV) & [Co(NH$_3$)$_4$(NO$_2$)$_2$]NO$_3$ & (iv) & sp$^3$d$^2$ & (S) & Optical\\
\hline
\end{tabular}
\end{center}
\end{bankpassage}
\begin{bankq}[id=SI12-P2-CHEM-S1-Q03, source={Score-I - Test 12 (09 Mar 2022)}, paper=2, qno=3, type=single, subtopic={Isomerism}, answer={A}, answer_src=key, difficulty=5, diff_parts={K2 S4 I0 C0 F0}, diff_note={Hybridisation, magnetism and isomerism type for four complexes, each option checked}, group=SI12-P2-CHEM-S1-G3-4]
Which is the only CORRECT combination?
\begin{qoptions}
\opt{A}{(I), (i), (Q)}
\opt{B}{(II), (iv), (P)}
\opt{C}{(I), (ii), (S)}
\opt{D}{(II), (iii), (Q)}
\end{qoptions}
\end{bankq}
\begin{bankq}[id=SI12-P2-CHEM-S1-Q04, source={Score-I - Test 12 (09 Mar 2022)}, paper=2, qno=4, type=single, subtopic={Isomerism}, answer={C}, answer_src=key, difficulty=5, diff_parts={K2 S4 I0 C0 F0}, diff_note={Same complex table; find the incorrect combination by checking isomerism and magnetism}, group=SI12-P2-CHEM-S1-G3-4]
Which is the only INCORRECT combination.
\begin{qoptions}
\opt{A}{(III), (ii), P}
\opt{B}{(IV), (ii), Q}
\opt{C}{(III), (ii), Q}
\opt{D}{(IV), (iii), Q}
\end{qoptions}

\noindent
\end{bankq}

\section*{TOAS papers}
\begin{bankq}[id=TOAS02-P1-CHEM-S1-Q06, source={TOAS - Test 02 (13 Feb 2022)}, paper=1, qno=6, type=single, subtopic={Isomerism}, answer={D}, answer_src=key, difficulty=6, diff_parts={K2 S4 I1 C0 F0}, diff_note={Count stereoisomers of four octahedral complexes including linkage isomers}]
The specie that show maximum number of stereo Isomers among the following.
\begin{qoptions}
\opt{A}{$\mathrm{[Ni(H_2O)_2(en)_2]^{2+}}$}
\opt{B}{$\mathrm{[Co(NH_3)_3(NO_2)_3]}$}
\opt{C}{$\mathrm{[Pt(py)_2(NH_3)_2Cl_2]^{2+}}$}
\opt{D}{$\mathrm{[CoBr_2Cl_2(NO_2)(SCN)]^{3-}}$}
\end{qoptions}
\end{bankq}
\begin{bankq}[id=TOAS02-P1-CHEM-S2-Q11, source={TOAS - Test 02 (13 Feb 2022)}, paper=1, qno=11, type=multi, subtopic={Isomerism}, answer={A,D}, answer_src=solved, working={Bidentate N,O ligand (not ambidentate); not planar overall (CH$_3$, H out of plane); stereoisomers cis-RR, cis-SS, cis-RS(meso), trans-RR, trans-SS, trans-RS(meso) = 6, with 2 enantiomeric pairs}, difficulty=7, diff_parts={K2 S4 I1 C0 F1}, diff_note={Stereoisomers of a square planar Pd bis(aminoacidate): cis/trans with chiral ligands}]
One of the isomer of $[\mathrm{Pd(NH_2{-}CH(CH_3){-}CO_2^-)_2}]$ is given as
\bfig{198.1}{TOAS02_p1_chem_s2_q11_1.png}
The correct statement(s) regarding the above compound is/are.
\begin{qoptions}
\opt{A}{Total stereoisomers will be 6}
\opt{B}{The ligand $\mathrm{(NH_2{-}CH(CH_3){-}COO^-)}$ is Ambidentate ligand}
\opt{C}{The molecule is planar since $\mathrm{Pd^{2+}}$ shows square planar geometry around it}
\opt{D}{Exactly two Enantiomeric pairs will be formed}
\end{qoptions}
\end{bankq}
\begin{bankq}[id=TOAS02-P1-CHEM-S3-Q05, source={TOAS - Test 02 (13 Feb 2022)}, paper=1, qno=5, type=numerical, subtopic={Nomenclature}, answer={4}, answer_src=solved, working={Incorrect: (a) ONO is O-bonded (nitrito-O), (b) ferrocene is $\eta^5$, (c) hydrogen hexachloridoplatinate(IV), (e) alphabetical order requires pentachloridonitrido; (d) and (f) correct}, difficulty=5, diff_parts={K1 S4 I0 C0 F1}, diff_note={Check names against IUPAC rules: bonding mode, hapticity, alphabetical order}]
Which of the following compounds/species are incorrectly matched against the indicated names.\\
(a) $\mathrm{[PdI_2(ONO)_2(H_2O)_2]}\to$ Diaquadiiodidodinitrito-Npalladium(IV)\\
(b) $\mathrm{Fe(\pi\text{-}C_5H_5)_2}\to$ Bis($\eta^1$--cyclopentadienyl)iron(II)\\
(c) $\mathrm{H_2[PtCl_6]}\to$ Hexachloridoplatinate(IV) hydride\\
(d) $\mathrm{[Al(OH)(H_2O)_5]SO_4}\to$ Pentaaquahydroxidoaluminium(III) sulphate\\
(e) $\mathrm{K_2[OsCl_5N]}\to$ Potassium nitridopentachloridoosmate(VI)\\
(f) $\mathrm{Cr(C_6H_6)_2}\to$ Bis($\eta$6--benzene)chromium(0)
\end{bankq}
\begin{bankq}[id=TOAS02-P1-CHEM-S3-Q06, source={TOAS - Test 02 (13 Feb 2022)}, paper=1, qno=6, type=numerical, subtopic={Nomenclature}, answer={20}, answer_src=solved, working={Secondary valencies: [Pd(NH$_3$)$_4$]Cl$_2$ 4, [Ni(NH$_3$)$_6$]Cl$_2$ 6, [Co(NH$_3$)$_4$Cl$_2$]Cl 6, [Pt(NH$_3$)$_2$Cl$_2$] 4; sum 20}, difficulty=4, diff_parts={K1 S3 I0 C0 F1}, diff_note={Free chloride from AgCl precipitate gives the coordination number of each complex}]
On the basis of the following observations made with aqueous solutions of given complexes, assigning secondary valency of each metal, write the sum of secondary valencies of all given complexes as your answer.
\par\nobreak
{\small\setlength{\tabcolsep}{3pt}\renewcommand{\arraystretch}{1.25}\begin{center}
\begin{tabular}{|L{2.9cm}|L{4.6cm}|}
\hline
Formula & Moles of AgCl precipitated per mole of compound with excess $\mathrm{AgNO_3}$ solution\\
\hline
(i)\ \ $\mathrm{PdCl_2.4NH_3}$ & 2 moles\\
\hline
(ii)\ \ $\mathrm{NiCl_2.6NH_3}$ & 2 moles\\
\hline
(iii)\ \ $\mathrm{CoCl_3.4NH_3}$ & 1 mole\\
\hline
(iv)\ \ $\mathrm{PtCl_2.2NH_3}$ & Zero mole\\
\hline
\end{tabular}\end{center}}
\end{bankq}
\begin{bankq}[id=TOAS04-P1-CHEM-S2-Q06, source={TOAS - Test 04 (15 Mar 2022)}, paper=1, qno=6, type=multi, subtopic={CFT}, answer={B,C}, answer_src=key, difficulty=6, diff_parts={K2 S3 I1 C0 F1}, diff_note={Spectrochemical series ranks splitting; absorbed wavelength is inversely related}]
\begin{tabular}[t]{@{}L{3.3cm}@{\hspace{0.8em}}L{3.2cm}@{}}
Co-ordination entity & Wave length of light absorbed (nm)\\
$\mathrm{[CoCl(NH_3)_5]^{2+}}$ & P\\
$\mathrm{[Co(NH_3)_5(H_2O)]^{3+}}$ & Q\\
$\mathrm{[Co(NH_3)_6]^{3+}}$ & R\\
$\mathrm{[Co(CN)_6]^{3-}}$ & S
\end{tabular}

Select the correct orders.
\begin{qoptions}
\opt{A}{P $>$ R $>$ Q}
\opt{B}{P $>$ Q $>$ R $>$ S}
\opt{C}{P $>$ Q $>$ R}
\opt{D}{Q $>$ P $>$ R $>$ S}
\end{qoptions}
\end{bankq}
\begin{bankpassage}[id=TOAS04-P1-CHEM-S3-G1-2, source={TOAS - Test 04 (15 Mar 2022)}, questions={1-2}]
Where two or more than two complexes having similar molecular weight (or) molecular formula but their structural arrangements are different these forms are known as isomer \& this phenomenon is known as isomerism; isomerism are of two types structural \& stereoisomerism ; stereoisomerism are of two types geometrical isomerism \& optical isomerism; compound which has plane of symmetry will not exhibit optical isomerism.
\end{bankpassage}
\begin{bankq}[id=TOAS04-P1-CHEM-S3-Q01, source={TOAS - Test 04 (15 Mar 2022)}, paper=1, qno=1, type=numerical, subtopic={Isomerism}, answer={10.00}, answer_src=key, difficulty=7, diff_parts={K2 S4 I1 C0 F1}, diff_note={Isomers of [Fe(CN)6]4- are linkage isomers with 0 to 6 isonitrile ligands, count geometrical forms}, group=TOAS04-P1-CHEM-S3-G1-2]
Total possible isomers of $\mathrm{K_4[Fe(CN)_6]}$.
\end{bankq}
\begin{bankq}[id=TOAS04-P1-CHEM-S3-Q02, source={TOAS - Test 04 (15 Mar 2022)}, paper=1, qno=2, type=numerical, subtopic={Isomerism}, answer={2.00}, answer_src=key, difficulty=5, diff_parts={K1 S3 I1 C0 F1}, diff_note={Geometrical isomers of an octahedral M(AA)(BB)(CC) complex}, group=TOAS04-P1-CHEM-S3-G1-2]
Total geometrical isomers theoretically possible in an octatedral complex of formula $\mathrm{[Co(C_2O_4)(dmg)(en)]}$ (In this complex $\mathrm{dmg^{-1}}$ is bonded by both `N' atoms with central metal)
\end{bankq}
\begin{bankq}[id=TOAS04-P2-CHEM-S1-Q01, source={TOAS - Test 04 (15 Mar 2022)}, paper=2, qno=1, type=multi, subtopic={CFT}, answer={A,B}, answer_src=key, difficulty=4, diff_parts={K1 S3 I0 C0 F1}, diff_note={Ligand field strength, metal and oxidation state effects on Delta_o}]
Which complex of the following pairs has the larger value of $\Delta_0$ ?\\
(i) $\mathrm{[Co(CN)_6]^{3-}}$ and $\mathrm{[Co(NH_3)_6]^{3+}}$\\
(ii) $\mathrm{[Co(H_2O)_6]^{3+}}$ and $\mathrm{[Rh(H_2O)_6]^{3+}}$\\
(iii) $\mathrm{[Co(H_2O)_6]^{3+}}$ and $\mathrm{[Co(H_2O)_3]^{3+}}$\\
(iv) $\mathrm{[Co(NH_3)_6]^{3+}}$ and $\mathrm{[CoF_6]^{3-}}$\\
Select the correct one
\begin{qoptions}
\opt{A}{$\mathrm{[Co(CN)_6]^{3-}>[Co(H_2O)_6]^{3+}}$}
\opt{B}{$\mathrm{[Co(H_2O)_6]^{2+}<[Co(H_2O)_6]^{3+}}$}
\opt{C}{$\mathrm{[Co(H_2O)_6]^{3+}>[Rh(H_2O)_6]^{3+}}$}
\opt{D}{$\mathrm{[Co(NH_3)_6]^{3+}<[CoF_6]^{3-}}$}
\end{qoptions}
\end{bankq}
\begin{bankq}[id=TOAS04-P2-CHEM-S1-Q02, source={TOAS - Test 04 (15 Mar 2022)}, paper=2, qno=2, type=multi, subtopic={Isomerism}, answer={B,D}, answer_src=key, difficulty=7, diff_parts={K2 S4 I1 C0 F1}, diff_note={Magnetic behaviour and type of isomerism for four pairs of complexes}]
Which of the following given pair of species exhibit similar kind of magnetic behaviour as well as at least one same kind of isomerism ?
\begin{qoptions}
\opt{A}{$\mathrm{[Zn(NH_3)_4]Cl_2}$ and $\mathrm{[Co(H_2O)_6]Cl_2}$}
\opt{B}{$\mathrm{[Cr(H_2O)_4Cl_2]^{+1}}$ and $\mathrm{[Co(NH_3)_3Cl_3]^{-1}}$}
\opt{C}{$\mathrm{K_3[Co(C_2O_4)_3]}$ and $\mathrm{K[Fe(EDTA)]}$ ; (EDTA = $\mathrm{EDTA^{4-}}$)}
\opt{D}{$\mathrm{[Co(NH_3)_5(NO_2)]SO_4}$ and $\mathrm{[Pd(PPh_3)_2(SCN)_2]}$}
\end{qoptions}
\end{bankq}
\begin{bankpassage}[id=TOAS04-P2-CHEM-S3-G1-2, source={TOAS - Test 04 (15 Mar 2022)}, questions={1-2}]
The metal--carbon bond in metal carbonyls posses both $\sigma$ \& $\pi$ character. The metal to ligand bonding creates a synergic effect which strengthens the bond between CO and the metal. Due to synergic effect M--C bond length \& C--O bond length affected.
\end{bankpassage}
\begin{bankq}[id=TOAS04-P2-CHEM-S3-Q01, source={TOAS - Test 04 (15 Mar 2022)}, paper=2, qno=1, type=numerical, subtopic={VBT}, answer={3}, answer_src=solved, working={(i) true: more back-donation in the electron-rich [Co(CO)$_4$]$^-$ lengthens C-O; (ii) false: weaker C-O bond in $x$; (iii) true; (iv) EAN of Co$^-$ with 4 CO and Mn$^+$ with 6 CO are both 36; so 3 statements}, difficulty=7, diff_parts={K2 S4 I1 C0 F1}, diff_note={Back bonding effect on C-O length and energy, M-C bond order and EAN for two carbonyls}, group=TOAS04-P2-CHEM-S3-G1-2]
Consider $\mathrm{[Co(CO)_4]^-}$ as $x$ and $\mathrm{[Mn(CO)_6]^+}$ as $y$ then find how many statements are correct?\\
(i) The C -- O bond length in $x$ is more than that in $y$\\
(ii) The C -- O bond energy in $x$ is more than that in $y$\\
(iii) Double bond character in metal to carbon bond in $x$ has more than that in $y$\\
(iv) Both $x$ and $y$ have the same EAN of the central metal atom
\end{bankq}
\begin{bankq}[id=TOAS04-P2-CHEM-S3-Q02, source={TOAS - Test 04 (15 Mar 2022)}, paper=2, qno=2, type=numerical, subtopic={VBT}, answer={8}, answer_src=solved, working={Longer than free CO: Ni(CO)$_4$, V(CO)$_6^-$, CO$_2$, H$_2$CO$_3$, CH$_3$COO$^-$, Ni(PMe$_3$)(CO)$_3$, Cr(CO)$_6$, Ni(PF$_3$)(CO)$_3$ ($X=8$); only CO$^+$ is shorter ($Y=1$); $X/Y=8$}, difficulty=7, diff_parts={K2 S4 I1 C0 F1}, diff_note={Count species with C-O length above or below free CO (back bonding, CO+, carbonates)}, group=TOAS04-P2-CHEM-S3-G1-2]
Consider the following species\\
$\mathrm{[Ni(CO)_4]}$, $\mathrm{[V(CO)_6]^-}$, $\mathrm{CO^+}$, $\mathrm{CO_2}$, $\mathrm{H_2CO_3}$, $\mathrm{CH_3COO^-}$, $\mathrm{[NiP(CH_3)_3(CO)_3]}$, $\mathrm{[Cr(CO)_6]}$, $\mathrm{[NiPF_3(CO)_3]}$\\
If bond length in CO is `$d$' in isolated molecule.\\
The number of species having higher C-O bond length then `$d$' is/are `$X$'\\
The number of species having lower C-O bond length then `$d$' is/are `$Y$'\\
Find the value of $\left(\dfrac XY\right)$
\end{bankq}
\begin{bankq}[id=TOAS04-P2-CHEM-S4-Q07, source={TOAS - Test 04 (15 Mar 2022)}, paper=2, qno=7, type=numerical, subtopic={CFT}, answer={36}, answer_src=key, difficulty=6, diff_parts={K2 S3 I1 C0 F1}, diff_note={d8 square planar: electrons in non-axial (6) and axial (2) d orbitals}]
If $x$ and $y$ are total number of electrons which are present in non-axial and axial set of d-orbitals respectively on central metal ion of $\mathrm{[Ni(DMG)_2]}$, then calculate value of $\dfrac{2x^2}{y}$ \underline{\hspace{1.5cm}}.
\end{bankq}
\begin{bankq}[id=TOAS04-P2-CHEM-S4-Q09, source={TOAS - Test 04 (15 Mar 2022)}, paper=2, qno=9, type=numerical, subtopic={Isomerism}, answer={7}, answer_src=key, difficulty=6, diff_parts={K2 S3 I1 C0 F1}, diff_note={Unpaired electrons 3, AgCl 1 mole, stereoisomers cis (2) and trans (1)}]
Consider the following complex for given three statements :\\
dichlorobis(ethylenediamine)chromium(III) chloride\\
(i) Number of unpaired electrons in complex $=(X)$\\
(ii) One mole of complex \bfigi{49.2}{TOAS04_p2_chem_s4_q09_1.png} $(Y)$ mole of AgCl.\\
(iii) Total stereo isomeric forms of above complex $=(Z)$ Find $X+Y+Z$
\end{bankq}
\begin{bankq}[id=TOAS06-P1-CHEM-S1-Q04, source={TOAS - Test 06 (17 Apr 2022)}, paper=1, qno=4, type=single, subtopic={Isomerism}, answer={C}, answer_src=key, difficulty=6, diff_parts={K2 S4 I1 C0 F0}, diff_note={Stereoisomer count for four coordination compounds in both geometries}]
Which of the following complexes has same number of stereoisomeric structures with any geometry (tetrahedral or square planar) ?
\begin{qoptions}
\opt{A}{$\mathrm{[M(NH_3)_2Cl_2]}$}
\opt{B}{$\mathrm{[M(NH_3)(py)(CN)(Cl)]}$}
\opt{C}{$\mathrm{[M(gly)_2]}$}
\opt{D}{$\mathrm{[M(H_2O)_2Br_2]}$}
\end{qoptions}
\end{bankq}
\begin{bankq}[id=TOAS06-P1-CHEM-S2-Q10, source={TOAS - Test 06 (17 Apr 2022)}, paper=1, qno=10, type=multi, subtopic={VBT}, answer={B,C}, answer_src=key, difficulty=6, diff_parts={K2 S3 I1 C0 F1}, diff_note={Pt(PPh3)4 is tetrahedral d10; XeO4 is tetrahedral: find the incorrect statements}]
Which of the following statement(s) is/are \textbf{INCORRECT}?
\begin{qoptions}
\opt{A}{$\mathrm{[Pt(PPh_3)_4]}$ is diamagnetic}
\opt{B}{$\mathrm{[Pt(PPh_3)_4]}$ is $dsp^2$ hybridised}
\opt{C}{$\mathrm{XeO_4}$ is a square planar}
\opt{D}{$\mathrm{XeO_4}$ and $\mathrm{[Pt(PPh_3)_4]}$ have same structure}
\end{qoptions}
\end{bankq}
\begin{bankq}[id=TOAS06-P1-CHEM-S2-Q12, source={TOAS - Test 06 (17 Apr 2022)}, paper=1, qno=12, type=multi, subtopic={Isomerism}, answer={A,B,C,D}, answer_src=key, difficulty=9, diff_parts={K2 S4 I2 C1 F1}, diff_note={Stereoisomers and geometrical isomers of [Co(gly)2(CN)Br]+ including linkage of CN}]
Compound $\mathrm{[Co(gly)_2(CN)(Br)]^+}$ ion has :
\begin{qoptions}
\opt{A}{only one optically inactive isomer}
\opt{B}{two linkage isomers}
\opt{C}{six geometrical isomers}
\opt{D}{eleven stereoisomers}
\end{qoptions}
\end{bankq}
\begin{bankq}[id=TOAS06-P1-CHEM-S3-Q05, source={TOAS - Test 06 (17 Apr 2022)}, paper=1, qno=5, type=numerical, subtopic={CFT}, answer={0.17}, answer_src=solved, working={Strong field CN$^-$: X Mn(II) $t_{2g}^5$ $a=-2.0$; Y Mn(III) $t_{2g}^4$ $b=-1.6$; Z Mn(I) $t_{2g}^6$ $c=-2.4$; $(a-b)/c=(-0.4)/(-2.4)=1/6\approx0.17$}, difficulty=5, diff_parts={K1 S3 I1 C0 F1}, diff_note={CFSE for Mn(II), Mn(III) and Mn(I) strong field cyanides}]
\bfig{143.7}{TOAS06_p1_chem_s3_q05_1.png}
if\\
the Crystal Field Stabilisation Energy of X is $=a*\Delta_o$\\
the Crystal Field Stabilisation Energy of Y is $=b*\Delta_o$\\
the Crystal Field Stabilisation Energy of Z is $=c*\Delta_o$\\
(Neglect pairing energy)\\
then the value of $(a-b)/c$ is
\end{bankq}
\begin{bankq}[id=TOAS07-P1-CHEM-S1-Q04, source={TOAS - Test 07 (01 May 2022)}, paper=1, qno=4, type=single, subtopic={VBT}, answer={D}, answer_src=key, difficulty=3, diff_parts={K1 S3 I0 C0 F0}, diff_note={EAN of carbonyls: Mn(CO)6 neutral gives 37}]
Which of the following does not obey Sidgwick effective atomic number rule -
\begin{qoptions}
\opt{A}{$\mathrm{Co_2(CO)_8}$}
\opt{B}{$\mathrm{Fe(CO)_5}$}
\opt{C}{$\mathrm{[V(CO)_6]^-}$}
\opt{D}{$\mathrm{[Mn(CO)_6]}$}
\end{qoptions}
\end{bankq}
\begin{bankq}[id=TOAS07-P1-CHEM-S2-Q09, source={TOAS - Test 07 (01 May 2022)}, paper=1, qno=9, type=multi, subtopic={CFT}, answer={A}, answer_src=key, difficulty=5, diff_parts={K1 S3 I1 C0 F1}, diff_note={Charge density and complex stability with EDTA: least tendency}]
$\mathrm{M^{2+}+NaH_2EDTA\to[M(EDTA)]^{2-}+2H^++2Na^+}$\\
Which of the following metal ion has least tendency for the above reaction.
\begin{qoptions}
\opt{A}{$\mathrm{Cs^+}$}
\opt{B}{$\mathrm{Mg^{2+}}$}
\opt{C}{$\mathrm{Ca^{2+}}$}
\opt{D}{$\mathrm{Pb^{2+}}$}
\end{qoptions}
\end{bankq}
\begin{bankq}[id=TOAS07-P1-CHEM-S2-Q10, source={TOAS - Test 07 (01 May 2022)}, paper=1, qno=10, type=multi, subtopic={Isomerism}, answer={A,D}, answer_src=key, difficulty=5, diff_parts={K1 S3 I1 C0 F1}, diff_note={Compare two hydrate/ionisation isomers: oxidation state, ligands, precipitate}]
Select similarity between : $\mathrm{[Co(NH_3)_4(H_2O)Cl]Cl_2}$ and $\mathrm{[Co(NH_3)_4Cl_2]Cl.H_2O}$
\begin{qoptions}
\opt{A}{Oxidation number and coordination number of metal}
\opt{B}{$\mathrm{Cl^-}$ and $\mathrm{H_2O}$ are the ligands}
\opt{C}{Amount of precipitate with $\mathrm{AgNO_3}$ solution (by equimolar 1 litre solution of both)}
\opt{D}{Number of Co--N linkages}
\end{qoptions}
\end{bankq}
\begin{bankq}[id=TOAS07-P1-CHEM-S4-Q09, source={TOAS - Test 07 (01 May 2022)}, paper=1, qno=9, type=numerical, subtopic={Nomenclature}, answer={2}, answer_src=key, difficulty=4, diff_parts={K1 S3 I0 C0 F1}, diff_note={Metal-carbon bonded complexes: olefin and carbonyl}]
Select complex in which metal--carbon linkage(s) is / are present :\\
$\left[\overset{\mathrm{III}}{\mathrm{Fe}}(\mathrm{EDTA})\right]^{-}$, $\mathrm{K[PtCl_3(\pi\text{--}C_2H_4)]}$, $\mathrm{[Cr(acac)_3]^\circ}$, $\mathrm{[Fe(CO)_5]}$, $\mathrm{[Co(C_2H_5O)_6]^{3-}}$, $\mathrm{[Ni(dmg)_2]}$
\end{bankq}
\begin{bankq}[id=TOAS07-P2-CHEM-S1-Q04, source={TOAS - Test 07 (01 May 2022)}, paper=2, qno=4, type=multi, subtopic={VBT}, answer={A,B}, answer_src=key, difficulty=6, diff_parts={K2 S3 I1 C0 F1}, diff_note={Back bonding lengthens C-O in more electron-rich carbonyls than Fe(CO)5}]
Which of the following metal carbonyl(s) have more C--O bond length than that of C--O bond in $\mathrm{Fe(CO)_5}$?
\begin{qoptions}
\opt{A}{$\mathrm{[V(CO)_6]^{\ominus}}$}
\opt{B}{$\mathrm{[Ti(CO)_6]^{2\ominus}}$}
\opt{C}{$\mathrm{[Mn(CO)_6]^{\oplus}}$}
\opt{D}{CO}
\end{qoptions}
\end{bankq}
\begin{bankq}[id=TOAS07-P2-CHEM-S1-Q05, source={TOAS - Test 07 (01 May 2022)}, paper=2, qno=5, type=multi, subtopic={Nomenclature}, answer={B,C,D}, answer_src=key, difficulty=6, diff_parts={K1 S4 I1 C0 F1}, diff_note={IUPAC names with oxidation states, ligand order and hapticity}]
Identify the correct option :
\begin{qoptions}
\opt{A}{$\mathrm{Na[PtBrCl(NH_3)(NO_2)]}$\newline IUPAC-Sodium amminebromidochloridonitrito-N-platinate(I)}
\opt{B}{$\mathrm{[Co(en)_3]Cl_3}$\newline IUPAC - tris(ethane-1,2-diamine)cobalt(III) chloride.}
\opt{C}{$\mathrm{[PtCl(NH_2Me)(NH_3)_2]Cl}$\newline IUPAC -diamminechlorido(methylamine)\newline platinum(II) chloride}
\opt{D}{$\mathrm{Fe(\pi\text{-}C_5H_5)_2}$\newline IUPAC-Bis($\eta^5$-cyclopentadienyl)iron(II)}
\end{qoptions}
\end{bankq}
\begin{bankq}[id=TOAS07-P2-CHEM-S1-Q06, source={TOAS - Test 07 (01 May 2022)}, paper=2, qno=6, type=multi, subtopic={Isomerism}, answer={B,C}, answer_src=key, difficulty=7, diff_parts={K2 S4 I1 C0 F1}, diff_note={Ionisation, hydrate, linkage and polymerisation isomerism statements}]
Choose correct statement/s.
\begin{qoptions}
\opt{A}{$\mathrm{[Co(NH_3)_5Br]SO_4}$ and $\mathrm{[Co(NH_3)_5SO_4]Br}$ both are ionization isomers of each other and in both complex all S-O bond length are identical}
\opt{B}{$\mathrm{[Cr(H_2O)_6]Cl_3}$ and $\mathrm{[Cr(H_2O)_4Cl_2]Cl.2H_2O}$ are hyrate isomers of each other and both isomers can be distinguished by conductivity measurements}
\opt{C}{$\mathrm{[Co(NH_3)_5NO_2]Cl_2}$ can show ionization isomerism as well as linkage isomerism}
\opt{D}{$\mathrm{[Pt(NH_3)_2Cl_2]}$ and $\mathrm{[Pt(NH_3)_4][PtCl_4]}$ are polymerization isomer of each other and both show electrical conductance in aqueous solution.}
\end{qoptions}
\end{bankq}
\begin{bankpassage}[id=TOAS07-P2-CHEM-S3-G5-6, source={TOAS - Test 07 (01 May 2022)}, questions={5-6}]
According to Werner's coordination theory every central metal atom show two type of valencies in the complexes and try to satisfy both of it's valencies.
\end{bankpassage}
\begin{bankq}[id=TOAS07-P2-CHEM-S3-Q05, source={TOAS - Test 07 (01 May 2022)}, paper=2, qno=5, type=numerical, subtopic={Nomenclature}, answer={0.33}, answer_src=key, difficulty=4, diff_parts={K1 S3 I0 C0 F1}, diff_note={Werner theory for CoCl3.5NH3: ions per molecule}, group=TOAS07-P2-CHEM-S3-G5-6]
For complex $\mathrm{CoCl_3.5NH_3}$\\
If P.V. $=x$\\
\hspace*{1em}S.V. $=y$\\
and total number of ions in aqueous solution per molecule of complex is $=z$\\
then calculate $\dfrac{z}{x+y}=$ ?
\end{bankq}
\begin{bankq}[id=TOAS07-P2-CHEM-S3-Q06, source={TOAS - Test 07 (01 May 2022)}, paper=2, qno=6, type=numerical, subtopic={Nomenclature}, answer={2.00}, answer_src=key, difficulty=5, diff_parts={K1 S3 I1 C0 F1}, diff_note={PtCl4.4NH3 has two ionisable chlorides}, group=TOAS07-P2-CHEM-S3-G5-6]
How many moles of AgCl(ppt.) is obtained when one mol aqueous solution of $\mathrm{PtCl_4.4NH_3}$ is treated with excess of $\mathrm{AgNO_3}$ solution.
\end{bankq}
\begin{bankq}[id=TOAS07-P2-CHEM-S4-Q08, source={TOAS - Test 07 (01 May 2022)}, paper=2, qno=8, type=numerical, subtopic={CFT}, answer={3}, answer_src=key, difficulty=5, diff_parts={K1 S3 I1 C0 F1}, diff_note={Ligands with both pi donor and pi* acceptor character}]
How many given ligands can act as $\pi$-electron donor and $\pi^*$-acceptor in complexes\\
$\mathrm{CH_2=CH_2}$, $\mathrm{NH_3}$, $\mathrm{PF_3}$, $\mathrm{C_6H_6}$, $\mathrm{C_5H_5^-}$, $\mathrm{H_2O}$
\end{bankq}
\begin{bankq}[id=TOAS07-P2-CHEM-S4-Q09, source={TOAS - Test 07 (01 May 2022)}, paper=2, qno=9, type=numerical, subtopic={VBT}, answer={0}, answer_src=key, difficulty=4, diff_parts={K1 S3 I0 C0 F1}, diff_note={EAN difference for Fe(CO)5 and Fe(CO)2(NO)2}]
What is the difference between the EAN value of $\mathrm{[Fe(CO)_5]}$ \& $\mathrm{[Fe(CO)_2(NO)_2]}$
\end{bankq}
\begin{bankq}[id=TOAS08-P1-CHEM-S1-Q01, source={TOAS - Test 08 (15 May 2022)}, paper=1, qno=1, type=multi, subtopic={Isomerism}, answer={D}, answer_src=key, difficulty=7, diff_parts={K2 S4 I1 C0 F1}, diff_note={Isomer counts, bridging carbonyl bond length, IUPAC name, orbital type for four complexes}]
Which of the following is(are) correct ?
\begin{qoptions}
\opt{A}{Number of geometrical isomers of ionisation isomers of $[\mathrm{Co(NH_3)_3(H_2O)Cl_2}]\mathrm{Br}$ is 5 (except the given compound)}
\opt{B}{In $\mathrm{Fe_2(CO)_9}$, $d_{\mathrm{C-O}}$ (bridging) $<$ $d_{\mathrm{C-O}}$ (terminal)}
\opt{C}{IUPAC name of $[(\mathrm{Cl_3Sn})_2\mathrm{RhCl_2Rh(SnCl_3)_2}]^{4-}$ is Di-$\mu$-chloridotetrakis(trichlorostannito)dirhodium(I) ion}
\opt{D}{Complexes $[\mathrm{Co(NO_2)_6}]^{4-}$ and $[\mathrm{CoF_6}]^{3-}$ are inner orbital complex and outer orbital complex respectively}
\end{qoptions}
\end{bankq}
\begin{bankpassage}[id=TOAS08-P1-CHEM-S2-G7-8, source={TOAS - Test 08 (15 May 2022)}, questions={7-8}]
The complex $\mathrm{HgCo.4SCN}$ has two isomers $X$ and $Y$ (having $\mathrm{Co}$ as central metal ion in both) which are HOMOLEPTIC and magnetic moment ($\mu$) of $X=Y$.
\end{bankpassage}
\begin{bankq}[id=TOAS08-P1-CHEM-S2-Q07, source={TOAS - Test 08 (15 May 2022)}, paper=1, qno=7, type=single, subtopic={Isomerism}, answer={C}, answer_src=key, difficulty=6, diff_parts={K2 S4 I1 C0 F0}, diff_note={Isomers X and Y of Hg[Co(SCN)4]: structures, oxidation states, magnetic moment}, group=TOAS08-P1-CHEM-S2-G7-8]
Choose the incorrect option for $X$ and $Y$?\newline
(I) $X$ and $Y$ are different in structure around central metal ion\newline
(II) Oxidation of Hg and Co are same in both the complexes\newline
(III) $X$ has spin only magnetic moment equals to 1.73 BM\newline
(IV) Both isomers having same hybridisation around central metal ion
\begin{qoptions}
\opt{A}{I and II}
\opt{B}{III and IV}
\opt{C}{I and III}
\opt{D}{II, III and IV}
\end{qoptions}
\end{bankq}
\begin{bankq}[id=TOAS08-P1-CHEM-S2-Q08, source={TOAS - Test 08 (15 May 2022)}, paper=1, qno=8, type=single, subtopic={Isomerism}, answer={C}, answer_src=key, difficulty=6, diff_parts={K2 S3 I2 C0 F0}, diff_note={Linkage isomers of the two forms}, group=TOAS08-P1-CHEM-S2-G7-8]
The individual sum of the linkage isomers of `X' and `Y' are.
\begin{qoptions}
\opt{A}{8}
\opt{B}{6}
\opt{C}{10}
\opt{D}{4}
\end{qoptions}
\end{bankq}
\begin{bankq}[id=TOAS08-P2-CHEM-S1-Q01, source={TOAS - Test 08 (15 May 2022)}, paper=2, qno=1, type=multi, subtopic={Isomerism}, answer={B,C,D}, answer_src=key, difficulty=7, diff_parts={K2 S4 I1 C0 F1}, diff_note={Identify Pt/Co complexes from clues; isomerism and unpaired electrons}]
Element ``$X$'' forms a compound $\mathrm{XCl_2.4NH_3}$ that does not show any type of isomerism. Element ``$Y$'' forms $\mathrm{YCl_3.4NH_3}$ type compound which has two isomers ``$P$'' \& ``$Q$''. Also, both elements ``$X$'' \& ``$Y$'' belong to VIII group of Mendeleev's periodic table.\newline
``$Y$'' is the only metal present in vitamin $\mathrm{B_{12}}$ and compounds of ``$X$'' inhibit growth of tumours.\newline
Which is/are correct about $X$ \& $Y$ or their compounds?
\begin{qoptions}
\opt{A}{If one $\mathrm{NH_3}$ is removed from the formula of $\mathrm{YCl_3.4NH_3}$, then one of the isomer of newly formed complex compound becomes optically active}
\opt{B}{Both ``$P$'' \& ``$Q$'' are $1:1$ electrolytes}
\opt{C}{If two $\mathrm{NH_3}$ in $\mathrm{XCl_2.4NH_3}$ get replaced by Cl then the formula of the new compound will show stereoisomerism}
\opt{D}{$X$ and $X^{2+}$ have same number of unpaired electrons}
\end{qoptions}
\end{bankq}
\begin{bankq}[id=TOAS08-P2-CHEM-S2-Q07, source={TOAS - Test 08 (15 May 2022)}, paper=2, qno=7, type=single, subtopic={Isomerism}, answer={A}, answer_src=key, difficulty=8, diff_parts={K2 S4 I2 C1 F0}, diff_note={Isomerism and geometry of four metal complexes}]
Match the List.
\begin{center}\footnotesize\renewcommand{\arraystretch}{1.1}\setlength{\extrarowheight}{3.5pt}\setlength{\tabcolsep}{2pt}
\begin{tabular}{|c|l|c|l|}
\hline
\multicolumn{2}{|c|}{\parbox[t]{2.8cm}{\centering List-I\\ Complex Compound/Ion}} & \multicolumn{2}{c|}{\parbox[t]{2.8cm}{\centering List-II\\ Characteristics}}\\[3pt]
\hline
(P) & $[\mathrm{Zn(gly)_2}]$ & (I) & \parbox[t]{2.6cm}{Exhibits Geometrical Isomerism only}\\[3pt]
(Q) & $[\mathrm{Pd(NH_2CH(CH_3)COO^-)_2}]$ & (II) & \parbox[t]{2.6cm}{Exhibits Optical isomerism only}\\[3pt]
(R) & $[\mathrm{Co(C_2S_2O_2)_2(NH_3)_2}]^-$ & (III) & \parbox[t]{2.6cm}{Both geometrical \& optical isomerism are exhibited}\\[3pt]
(S) & $[\mathrm{Cu(NH_3)_4}][\mathrm{CuCl_4}]$ & (IV) & \parbox[t]{2.6cm}{Structural isomerism is exhibited}\\[3pt]
 & & (V) & \parbox[t]{2.6cm}{Square planar around any one central atom}\\[3pt]
 & & (VI) & \parbox[t]{2.6cm}{Tetrahedral around any one central atom}\\[3pt]
\hline
\end{tabular}
\end{center}
Which of the following option is correct match?
\begin{qoptions}
\opt{A}{$P\to VI;Q\to III;R\to III;S\to VI$}
\opt{B}{$P\to IV;Q\to V;R\to IV;S\to V$}
\opt{C}{$P\to V;Q\to II;R\to VI;S\to I$}
\opt{D}{$P\to II;Q\to I;R\to I;S\to III$}
\end{qoptions}
\end{bankq}
\begin{bankpassage}[id=TOAS09-P1-CHEM-S3-G1-2, source={TOAS - Test 09 (30 May 2022)}, questions={1-2}]
The complex compounds show space isomerism when the ligand species have different spatial arrangement surrounding the CMA
\end{bankpassage}
\begin{bankq}[id=TOAS09-P1-CHEM-S3-Q01, source={TOAS - Test 09 (30 May 2022)}, paper=1, qno=1, type=numerical, subtopic={Isomerism}, answer={3.00}, answer_src=key, difficulty=7, diff_parts={K2 S3 I1 C1 F1}, diff_note={Isomerism: standard approach for this subtopic; 2 concept(s), 3 step(s), needs a non-obvious insight, heavier computation}, group=TOAS09-P1-CHEM-S3-G1-2]
What is the number of stereo isomeric forms for the complex [Pt (py) (NH$_3$) (Br) (I)]
\end{bankq}
\begin{bankq}[id=TOAS09-P1-CHEM-S3-Q02, source={TOAS - Test 09 (30 May 2022)}, paper=1, qno=2, type=numerical, subtopic={Isomerism}, answer={3.00}, answer_src=key, difficulty=7, diff_parts={K2 S3 I1 C1 F1}, diff_note={Isomerism: standard approach for this subtopic; 2 concept(s), 3 step(s), needs a non-obvious insight, heavier computation}, group=TOAS09-P1-CHEM-S3-G1-2]
How many stereo isomers are possible for [Cr(H$_2$O)$_2$ (C$_2$O$_4$)$_2$]$^{-}$ ?
\end{bankq}
\begin{bankq}[id=TOAS09-P1-CHEM-S4-Q07, source={TOAS - Test 09 (30 May 2022)}, paper=1, qno=7, type=numerical, subtopic={CFT}, answer={5}, answer_src=key, difficulty=8, diff_parts={K2 S3 I2 C1 F1}, diff_note={CFT: standard approach for this subtopic; 2 concept(s), 3 step(s), needs a non-obvious insight, heavier computation}]
Total number of complex compound which have low spin as well as Diamagnetic.\par
[Fe(CN)$_6$]$^{3-}$, [Co(H$_2$O)$_6$]$^{3\oplus}$, [NiF$_6$]$^{2-}$, [Cr(H$_2$O)$_6$]$^{3\oplus}$, [Fe(CN)$_6$]$^{4-}$, [Co(C$_2$O$_4$)$_3$]$^{3-}$, [Ni(en)$_3$]$^{3+}$, [PtCl$_4$]$^{2-}$
\end{bankq}
\begin{bankq}[id=TOAS09-P2-CHEM-S1-Q03, source={TOAS - Test 09 (30 May 2022)}, paper=2, qno=3, type=single, subtopic={Nomenclature}, answer={A}, answer_src=key, difficulty=1, diff_parts={K1 S1 I0 C0 F0}, diff_note={Nomenclature: standard approach for this subtopic; 1 concept(s), 1 step(s)}]
Select the incorrect IUPAC name from the following
\begin{qoptions}
\opt{A}{H[AuCl$_4$] : Hydrogen tetrachloridoaurate(III)}
\opt{B}{[PtCl$_6$]$^{2-}$ : Hexachloridoplatinate(IV) ion}
\opt{C}{[Cu(NH$_3$)$_4$]SO$_4$ : Tetraamminecopper(II) sulphate}
\opt{D}{[Cr(en)$_3$]Cl$_3$ : Tris (ethylenediamine)chromium(III) chloride}
\end{qoptions}
\end{bankq}
\begin{bankq}[id=TOAS09-P2-CHEM-S1-Q04, source={TOAS - Test 09 (30 May 2022)}, paper=2, qno=4, type=single, subtopic={VBT}, answer={D}, answer_src=key, difficulty=2, diff_parts={K1 S2 I0 C0 F0}, diff_note={VBT: standard approach for this subtopic; 1 concept(s), 2 step(s)}]
Spin only magnetic moment in BM of [Fe(CO)$_4$(C$_2$O$_4$)]$^+$ is :
\begin{qoptions}
\opt{A}{5.92}
\opt{B}{0}
\opt{C}{2.82}
\opt{D}{1.73}
\end{qoptions}
\end{bankq}
\begin{bankq}[id=TOAS09-P2-CHEM-S1-Q05, source={TOAS - Test 09 (30 May 2022)}, paper=2, qno=5, type=single, subtopic={Isomerism}, answer={C}, answer_src=key, difficulty=1, diff_parts={K1 S1 I0 C0 F0}, diff_note={Isomerism: standard approach for this subtopic; 1 concept(s), 1 step(s)}]
Which of the following complex is optically active :
\begin{qoptions}
\opt{A}{trans-[Co(NH$_3$)$_4$Cl$_2$]$^+$}
\opt{B}{[Cr(H$_2$O)$_6$]$^{3+}$}
\opt{C}{cis-[Co(NH$_3$)$_2$(en)$_2$]$^{3+}$}
\opt{D}{trans-[Co(NH$_3$)$_2$(en)$_2$]$^{3+}$}
\end{qoptions}
\end{bankq}
\begin{bankq}[id=TOAS09-P2-CHEM-S2-Q03, source={TOAS - Test 09 (30 May 2022)}, paper=2, qno=3, type=numerical, subtopic={Isomerism}, answer={3}, answer_src=solved, working={Ionisation (Br$^-$/NO$_3^-$), hydrate (water inside/outside) and geometrical isomerism; no ambidentate ligand or ligand isomerism}, difficulty=7, diff_parts={K2 S3 I1 C1 F1}, diff_note={Isomerism: standard approach for this subtopic; 2 concept(s), 3 step(s), needs a non-obvious insight, heavier computation}]
Find the total number of type of isomerism (from given below) shown by the complex\par
\textbf{[Fe(NH$_3$)$_4$(H$_2$O)Br] (NO$_3$)$_2$}\par
Ionization isomerism, linkage isomerism, ligand isomerism, Hydrate isomerism, geometrical isomerism
\end{bankq}
\begin{bankq}[id=TOAS09-P2-CHEM-S2-Q04, source={TOAS - Test 09 (30 May 2022)}, paper=2, qno=4, type=numerical, subtopic={Isomerism}, answer={3}, answer_src=solved, working={Square planar Mabcd: 3 geometrical isomers}, difficulty=7, diff_parts={K2 S3 I1 C1 F1}, diff_note={Isomerism: standard approach for this subtopic; 2 concept(s), 3 step(s), needs a non-obvious insight, heavier computation}]
Total number of geometrical isomers for the complex [RhCl(CO)(PPh$_3$) (NH$_3$)] is:
\end{bankq}
\end{document}
